\chapter{Theoretical Background and Literature Review}

% Chapter Introduction
This chapter provides a comprehensive review of the theoretical foundations and current state-of-the-art in Retrieval-Augmented Generation (RAG) systems, with particular emphasis on query enhancement techniques and their application to Arabic language processing. The chapter is structured to guide the reader from fundamental RAG concepts to the specific challenges of Arabic language retrieval, establishing the theoretical foundation for our proposed methodology.

\section{Theoretical Background: RAG Fundamentals}

\subsection{The Hallucination Problem in Large Language Models}

Large Language Models (LLMs) have demonstrated remarkable capabilities in natural language understanding and generation. However, they suffer from a fundamental limitation: their knowledge is parametric, encoded in model weights during training, making it static and prone to hallucinations \cite{song_2024_a}. When queried about information beyond their training data or requiring up-to-date knowledge, LLMs may generate plausible but factually incorrect responses.

This limitation becomes particularly problematic in domain-specific applications where accuracy is critical. The parametric nature of LLM knowledge means that:
\begin{enumerate}
    \item Knowledge becomes outdated as time passes since training
    \item Domain-specific or specialized information may be underrepresented in training data
    \item The model cannot distinguish between what it knows with certainty and what it is inferring
    \item Updating knowledge requires expensive retraining of the entire model
\end{enumerate}

\subsection{The Standard RAG Architecture}

Retrieval-Augmented Generation (RAG) addresses the hallucination problem by augmenting LLMs with external knowledge retrieval \cite{lewis_2021_retrievalaugmented}. The standard RAG architecture consists of two primary components:

\textbf{The Retriever:} Responsible for finding relevant documents or passages from a knowledge base given a user query. Retrieval methods can be categorized as:
\begin{itemize}
    \item \textbf{Sparse Retrieval:} Traditional methods like BM25 that rely on exact term matching and statistical term weighting
    \item \textbf{Dense Retrieval:} Neural methods that encode queries and documents into dense vector representations, enabling semantic similarity matching \cite{karpukhin_2020_dense}
    \item \textbf{Hybrid Retrieval:} Combinations of sparse and dense methods to leverage complementary strengths
\end{itemize}

\textbf{The Generator:} An LLM that synthesizes the final answer by conditioning on both the user query and the retrieved context. The retrieved passages provide grounding, reducing hallucinations and enabling the model to cite sources.

The standard RAG pipeline operates as follows:
\begin{enumerate}
    \item User submits a query $q$
    \item Retriever finds top-$k$ relevant passages $\{d_1, d_2, ..., d_k\}$ from knowledge base $\mathcal{D}$
    \item Generator produces answer $a$ conditioned on $q$ and retrieved passages
    \item System returns answer $a$ (optionally with source citations)
\end{enumerate}

\section{The Retrieval Bottleneck and Query Mismatch}

\subsection{The Semantic Gap in Retrieval}

Despite advances in dense retrieval methods, a fundamental challenge persists: the semantic gap between user queries and relevant documents. This gap manifests when queries and documents use different vocabulary, phrasing, or levels of abstraction to express related concepts. Vector-based semantic search, while more robust than keyword matching, still struggles when:

\begin{itemize}
    \item Queries are too short or ambiguous to capture user intent
    \item Users employ colloquial language while documents use formal terminology
    \item Queries contain domain-specific jargon not present in the knowledge base
    \item Morphological variations create lexical mismatches
\end{itemize}

This semantic gap directly impacts retrieval recall—the proportion of relevant documents successfully retrieved. Low recall in the retrieval stage creates a bottleneck that cannot be compensated for by the generator, as the LLM can only work with the passages it receives.

\subsection{Query Sensitivity and Noise}

Recent research has highlighted the sensitivity of RAG systems to query formulation. Studies on RAG robustness demonstrate that minor variations in query phrasing can significantly impact retrieval quality \cite{perin_2025_investigating}. Query noise can originate from multiple sources:

\begin{itemize}
    \item \textbf{Typographical errors:} Spelling mistakes and typos that prevent exact matching
    \item \textbf{Linguistic variations:} Dialectal differences, informal language, and code-switching
    \item \textbf{Morphological complexity:} Rich inflectional systems that create numerous surface forms
    \item \textbf{Incomplete queries:} Users providing insufficient context or using ambiguous terms
\end{itemize}

The QE-RAG framework specifically addresses query noise, demonstrating that query enhancement techniques can significantly improve retrieval robustness \cite{zhang_2025_qerag}. Crucially, the concept of "noise" in English queries (primarily typos and paraphrasing) is analogous to dialectal variations in Arabic—both represent deviations from the standard form present in the knowledge base.

\section{Taxonomy of Advanced RAG Paradigms}

Beyond the standard RAG architecture, researchers have proposed various enhancements targeting different components of the pipeline. We categorize these approaches into three main paradigms:

\subsection{Index-Centric Approaches: Structural Enhancements}

Index-centric methods modify how knowledge is stored and organized to improve retrieval effectiveness:

\textbf{Hierarchical Structures:} RAPTOR \cite{sarthi_2024_raptor} introduces recursive abstractive processing, organizing documents into tree structures where higher levels contain summaries of lower levels. This enables retrieval at multiple levels of abstraction.

\textbf{Graph-Based Indexing:} LightRAG \cite{guo_2024_lightrag} and similar approaches represent knowledge as graphs, capturing relationships between entities and concepts. Graph structures enable multi-hop reasoning and relationship-aware retrieval \cite{han_2024_retrievalaugmented}.

\textbf{Proposition-Based Chunking:} Dense X Retrieval \cite{chen_2024_dense} decomposes documents into atomic propositions, creating fine-grained retrieval units that reduce context dilution.

While these approaches show promise, they primarily address the storage and indexing side of RAG, requiring significant preprocessing and potentially complex query routing mechanisms.

\subsection{Process-Centric Approaches: Agentic RAG}

Process-centric methods introduce reasoning loops and decision-making into the RAG pipeline:

\textbf{Iterative Retrieval:} Systems that perform multiple retrieval rounds, refining queries based on initial results or intermediate generation steps.

\textbf{Self-Reflective RAG:} Approaches like Self-RAG \cite{asai_2024_selfrag} that incorporate critique tokens, enabling the model to assess retrieval quality and decide when additional retrieval is needed.

\textbf{Agentic Workflows:} Complex systems that decompose queries, plan retrieval strategies, and orchestrate multiple tools and knowledge sources.

These approaches offer flexibility and can handle complex queries, but introduce additional latency and complexity, making them challenging to optimize and deploy.

\section{Query Enhancement Techniques: State-of-the-Art}

Query enhancement represents a third paradigm: modifying the input query before retrieval to improve matching with relevant documents. This approach is attractive because it operates as a modular layer that can be added to existing RAG systems without modifying the index or generator.

\subsection{Generative Query Expansion}

Generative query expansion leverages LLMs to create enhanced or alternative queries:

\textbf{HyDE (Hypothetical Document Embeddings):} Proposed by Gao et al. \cite{gao_2022_precise}, HyDE instructs an LLM to generate a hypothetical document that would answer the query, then uses this generated document for retrieval instead of the original query. The intuition is that document-to-document similarity is more reliable than query-to-document similarity. HyDE operates in a zero-shot manner, requiring no training data, and has shown effectiveness for both sparse (BM25) and dense retrieval methods.

\textbf{Query2Doc:} Extends the HyDE concept by generating multiple pseudo-documents and combining them with the original query, creating a richer representation for retrieval \cite{wang_2023_query2doc}.

The key trade-off in generative expansion is improved recall versus increased latency and computational cost, as LLM inference is required for each query.

\subsection{Query Rewriting and Refinement}

Query rewriting transforms the original query into a more effective form:

\textbf{Rewrite-Retrieve-Read:} A framework that explicitly separates query rewriting from retrieval, using an LLM to reformulate queries before searching \cite{ma_2023_query}.

\textbf{Query Optimization:} Systematic approaches to query transformation, including expansion with synonyms, simplification of complex queries, and normalization of linguistic variations \cite{song_2024_a, louis_2024_query}.

\textbf{Dialect Normalization:} In multilingual contexts, query rewriting can translate dialectal or colloquial queries into the standard form present in the knowledge base.

\subsection{Iterative and Active Retrieval}

Advanced query enhancement incorporates feedback loops:

\textbf{RQ-RAG (Refining Queries for RAG):} Introduces an "active generator" architecture where the LLM iteratively refines queries based on retrieval results, learning when to stop searching \cite{chan_2024_rqrag}. This approach treats query refinement as a reinforcement learning problem.

\textbf{Self-RAG:} Incorporates special tokens that enable the model to critique retrieval quality and decide whether to retrieve additional information \cite{asai_2024_selfrag}.

\textbf{LevelRAG:} Implements a two-tiered architecture with high-level planning (query decomposition, tool selection) and low-level execution (sparse, dense, and web search), demonstrating the value of orchestrating multiple retrieval strategies \cite{wang_2025_levelrag}.

\section{RAG in Morphologically Rich and Low-Resource Languages}

\subsection{The Arabic Linguistic Challenge}

Arabic presents unique challenges for RAG systems due to its linguistic characteristics:

\textbf{Morphological Richness:} Arabic employs a root-and-pattern morphology where words are derived from tri-consonantal roots through templatic patterns and affixation. This creates numerous surface forms for related concepts, complicating lexical matching.

\textbf{Diglossia:} Arabic exhibits a diglossic situation with Modern Standard Arabic (MSA) used in formal writing and various regional dialects used in speech and informal communication. Dialectal Arabic differs significantly from MSA in vocabulary, morphology, and syntax.

\textbf{Orthographic Variations:} Arabic script allows variations in representing certain phonemes (e.g., Hamza: ء، أ، إ، ؤ، ئ; Ya: ي، ى; Alif: ا، آ), leading to spelling inconsistencies even within MSA.

\textbf{Diacritical Marks:} While diacritics (tashkeel) disambiguate pronunciation and meaning, they are typically omitted in modern texts, increasing ambiguity.

These characteristics create a significant "morphological and dialectal gap" between user queries (often dialectal, informal, with spelling variations) and knowledge base documents (typically MSA, formal, standardized).

\subsection{Current Arabic RAG Baselines}

Recent work has begun establishing baselines for Arabic RAG systems:

\textbf{Embedding Models:} Multilingual embedding models like BGE-m3 have shown reasonable performance on Arabic retrieval tasks. The MIRACL benchmark \cite{zhang_2023_miracl} provides standardized evaluation for Arabic information retrieval.

\textbf{Chunking Strategies:} Research by Alsubhi et al. \cite{alsubhi_2025_optimizing} demonstrates that sentence-aware chunking outperforms fixed-size chunking for Arabic, suggesting that linguistic structure matters for retrieval quality. Additional work by El-Beltagy and Abdallah \cite{elbeltagysamhaar_2024_exploring} has explored RAG applications in Arabic contexts.

\textbf{Hybrid Approaches:} Combining BM25 with dense retrieval has proven effective for Arabic, leveraging the complementary strengths of keyword matching and semantic similarity.

However, these baselines primarily focus on MSA and do not explicitly address dialectal variations or morphological complexity in query formulation.

\subsection{The Research Gap}

While substantial progress has been made in two areas—query enhancement techniques for English (Section 2.4) and baseline RAG systems for Arabic (Section 2.5.2)—the intersection remains under-explored. Specifically:

\begin{enumerate}
    \item \textbf{Query Enhancement for Arabic:} Most query enhancement research focuses on English, where the primary "noise" is typos and paraphrasing. The applicability of techniques like HyDE and query rewriting to Arabic's morphological and dialectal challenges is unclear.
    
    \item \textbf{Dialectal Query Handling:} Existing Arabic RAG systems assume MSA queries. Real-world users, however, often formulate queries in dialectal Arabic or mix dialects with MSA (code-switching). How to bridge this dialect-document gap through query enhancement is an open question.
    
    \item \textbf{Evaluation Frameworks:} While MIRACL provides a benchmark for Arabic retrieval, it focuses on MSA queries. Evaluating query enhancement techniques for dialectal or morphologically complex queries requires additional datasets or synthetic data generation strategies.
    
    \item \textbf{Technique Adaptation:} It is unclear whether English query enhancement techniques transfer directly to Arabic or require language-specific adaptations. For example, does HyDE's hypothetical document generation work effectively with Arabic LLMs? Does query rewriting need to incorporate morphological normalization?
\end{enumerate}

\section{Summary and Research Positioning}

This chapter has established the theoretical foundation for our research:

\begin{itemize}
    \item RAG systems address LLM hallucinations by incorporating external knowledge retrieval (Section 2.1)
    \item Retrieval quality is bottlenecked by query-document mismatch and query noise (Section 2.2)
    \item While various RAG enhancements exist (Section 2.3), query enhancement offers a modular, effective approach (Section 2.4)
    \item Arabic presents unique challenges due to morphology and diglossia (Section 2.5.1)
    \item Current Arabic RAG baselines exist but do not address dialectal query enhancement (Section 2.5.2)
    \item A research gap exists at the intersection of query enhancement techniques and Arabic linguistic challenges (Section 2.5.3)
\end{itemize}

Our research addresses this gap by investigating how query enhancement techniques can be adapted and applied to improve retrieval recall in Arabic RAG systems, with particular attention to the morphological and dialectal challenges that characterize real-world Arabic queries. The following chapter presents our methodology for systematically evaluating query enhancement approaches in the Arabic context.
